\chapter{Gate-Sizing Optimization}
\label{ch:optimization}

The optimizer is implemented in
\code{optimization/optimizer.py:CircuitOptimizer}, with a deliberately
narrow public surface:

\begin{lstlisting}[style=pythonstyle, numbers=none]
class CircuitOptimizer:
    def __init__(self, circuit: Circuit,
                 heuristic: OptimizationHeuristic = BRUTE_FORCE,
                 random_seed: int | None = None) -> None: ...
    def optimize(self) -> OptimizationResult: ...
\end{lstlisting}

Everything else --- candidate generation, cost evaluation, the acceptance
policy, and termination --- is private to this class or to
\code{optimization/heuristics.py}, and is documented in this chapter by
name because the \emph{names} of these private functions are themselves
the clearest statement of the algorithm's structure.

\section{Cost Function}
\label{sec:opt-cost}

Every candidate gate-sizing change is scored by a single normalized cost
function that trades off circuit delay $D_{\mathrm{circuit}}$, total power
$P_{\mathrm{total}}$, total area $A_{\mathrm{total}}$, and any remaining
timing violation, each normalized against the corresponding value of the
\emph{original, unmodified} design (subscript $\mathrm{ref}$), so that the
cost is dimensionless and comparable across circuits of very different
absolute size:

\begin{equation}
    J = w_D \frac{D_{\mathrm{circuit}}}{D_{\mathrm{ref}}}
      + w_P \frac{P_{\mathrm{total}}}{P_{\mathrm{ref}}}
      + w_A \frac{A_{\mathrm{total}}}{A_{\mathrm{ref}}}
      + \lambda\, \frac{\max(0, -\mathrm{WNS})}{T_{\mathrm{ref}}}.
    \label{eq:cost}
\end{equation}

The weights $w_D, w_P, w_A, \lambda$ are read from
\code{Config.optimization.weights} (\code{Optimizati\newline onWeights},
\cref{tab:arch-domain-api}). $D_{\mathrm{ref}}, P_{\mathrm{ref}},
A_{\mathrm{ref}}$ are captured exactly once by the static method
\code{CircuitOptimizer.\_normalization\_reference(circuit, timing) ->
\_Normalization Reference}, called only from \code{\_initial\_state()} at
the very start of \code{optimize()}, and never recomputed thereafter, so
that $J$ always measures cost \emph{relative to the starting point}.
$T_{\mathrm{ref}}$, a normalizing time reference for the violation penalty,
is taken as the tightest required time among the circuit's primary outputs.

Cost itself is computed by the private method
\code{CircuitOptimizer.\_evaluate(circuit, reference, timing=None) ->
\_Evaluation}, where \code{\_Evaluation} is a frozen dataclass bundling
\code{circuit}, \code{timing}, and \code{cost} together --- the single
value threaded through every other private method below.

\section{Candidate Evaluation}
\label{sec:opt-eval}

Consistent with the immutability principle stated in
\cref{sec:arch-principles}, evaluating a candidate resize never mutates the
circuit currently being optimized.
\code{CircuitOptimizer.\_candidate\_ evaluation(current, reference, gate,
candidate\_cell) -> \_CandidateEvaluation | None} is the single choke
point through which every candidate --- from every heuristic --- passes,
in three ordered steps.

\begin{enumerate}
    \item \textbf{Cheap pre-filter.}
          \code{replacement\_area\_and\_power(current.circuit, gate.name,
          candidate\_cell)} (\cref{tab:arch-domain-api}) computes the
          resulting total area and power \emph{without} constructing a new
          \code{Circuit} object at all. If this alone already exceeds the
          configured budget
          (\code{\_violates\_area\_or\_power\_constraints}), the candidate
          is discarded immediately, before paying for a full circuit
          construction or a timing re-run.
    \item \textbf{Full construction and re-analysis.} Only if the cheap
          check passes is \code{replace\_gate \_cell(current.circuit,
          gate.name, candidate\_cell)} called to build the new, independent
          \code{Circuit} variant, which is then fully re-evaluated:
          \code{\_analyze(circuit)} runs \code{analyze\_timing}
          (incrementing an internal \code{self.\_sta\_calls} counter,
          reported as \code{OptimizationResult.sta\_calls}), and
          \code{\_satisfies\_constraints(circuit)} performs the exact
          (non-delta) area/power check against the fully rebuilt circuit.
    \item \textbf{Cost.} \code{\_evaluate(candidate\_circuit, reference)}
          scores the surviving candidate via \cref{eq:cost}.
\end{enumerate}

If a candidate is not selected, the \code{Circuit} object built for it is
simply discarded --- there is no revert step, because nothing was ever
mutated in place. This is functionally equivalent to an
evaluate-then-revert pattern built on a mutable graph, but structurally
eliminates the possibility of a forgotten or incomplete revert silently
corrupting a subsequent evaluation.

\section{The Acceptance Policy}
\label{sec:opt-accept}

A detail that is easy to miss without reading the source directly, but is
central to how the four heuristics relate to one another: \textbf{there is
exactly one implementation of the two-phase accept/reject policy, and every
heuristic calls it.}

\vspace{10pt}
\begin{lstlisting}[style=pythonstyle, numbers=none]
def _select_next_change(self, current: _Evaluation,
                        candidates: list[_CandidateEvaluation]
                        ) -> _AcceptedChange | _StopDecision: ...
@staticmethod
def _best_timing_change(current, candidates) -> _CandidateEvaluation | None: ...
@staticmethod
def _best_cost_change(candidates) -> _CandidateEvaluation | None: ...
\end{lstlisting}


\newpage \noindent \code{\_select\_next\_change} implements the two-phase rule directly:

\begin{enumerate}
    \item \textbf{While a timing violation remains}
          ($\mathrm{WNS} < 0$, \code{OptimizationPhase.TIMING\_REPAIR}):
          \code{\_best\_timing\_change} selects, among the supplied
          candidates, the one with the lowest cost $J$ among those that
          strictly improve WNS, or, failing that, strictly improve TNS with
          WNS unchanged. If none qualify, \code{\_select\_next\_change}
          returns a \code{\_StopDecision} with termination reason
          \code{NO\_TIMING\_IMPROVEMENT}.
    \item \textbf{Once timing is compliant}
          ($\mathrm{WNS} \geq 0$, \code{OptimizationPhase.COST\_REDUCTION}):
          \code{\_best\_cost\_change} selects the lowest-cost candidate that
          does not reintroduce a violation, and the change is accepted only
          if the resulting cost improvement exceeds
          \code{config.optimization.minimum\_cost\_improvement}. Otherwise, termination reason
          \code{TIMING\_MET\_NO\_COST\_IMPROVEMENT} is returned.
\end{enumerate}

What differs between heuristics is only \emph{how many} candidates
\code{\_select\_next\_change} is called with, and \emph{in what order}
they were generated (\cref{sec:opt-heuristics}) --- not the acceptance
rule itself. Brute force calls it once per iteration with the
\emph{entire} evaluated candidate batch:

\begin{lstlisting}[style=pythonstyle, numbers=none]
choices = tuple(choice_iterator)    # exhaust the whole iterator
candidates = self._candidate_evaluations(current, reference, choices)
return self._select_next_change(current, candidates)  # pick the best of all
\end{lstlisting}

while the other three heuristics instead call
\code{\_evaluate\_single\_choice}, which evaluates exactly \emph{one}
candidate and delegates to the very same function with a singleton list:

\begin{lstlisting}[style=pythonstyle, numbers=none]
def _evaluate_single_choice(self, current, reference, gate, candidate_cell
        ) -> _AcceptedChange | _RejectedChange:
    candidate = self._candidate_evaluation(current, reference, gate, candidate_cell)
    if candidate is None:
        return _RejectedChange(gate.name, candidate_cell.name,
                                "candidate violates area or power constraints")
    decision = self._select_next_change(current, [candidate]) 
    ...
\end{lstlisting}

This makes slack-weighted capacitance, criticality/effort-gap, and random
greedy \emph{sequential, first-improvement} searches: each pulls one
candidate at a time from a pre-ordered stream and immediately accepts or
rejects it, moving to the next candidate in that same stream on rejection,
without re-ranking until a change is actually accepted. Brute force, in
contrast, is a \emph{best-of-batch} search: every eligible candidate is
evaluated fresh each iteration, and the single best one is chosen. The
optimizer's outer loop, \code{optimize()}, regenerates the candidate stream
(\code{choice\_iterator = None}, forcing a fresh call to
\code{\_choice\_iterator} on the next iteration) only \emph{after} an
accepted change, so a rejected candidate does not trigger re-ranking ---
the search simply continues down the same, still-valid-for-the-current-circuit
ordering.

\section{Termination}
\label{sec:opt-termination}

\code{OptimizationTermination} enumerates five possible outcomes, recorded
verbatim in \code{Optim\newline izationResult.termination} and in the
\code{optimization\_*.csv} artifacts: \code{DISABLED} (optimization turned
off in configuration); \code{NO\_TIMING\_IMPROVEMENT} (still violating, no
candidate improves WNS/TNS); \code{NO\_FEASIBLE\_CHANGE} (the candidate
stream is exhausted --- only reachable by the three sequential heuristics,
since brute force always sees its full batch every iteration);
\code{TIMING\_MET\_NO\_COST\_IMPROVEMENT} (compliant, but no candidate
reduces cost enough); and \code{MAXIMUM\_ITERATIONS}.

\section{The Four Sizing Heuristics}
\label{sec:opt-heuristics}

\code{OptimizationHeuristic} (\code{optimization/heuristics.py}) has four
members: \code{BRUTE\_FORCE}, \code{SLACK\_WEIGHTED\_CAPACITANCE},
\code{CRITICALITY\_EFFORT\_GAP}, and \code{RANDOM\_GREEDY}. All four are
dispatched from a single method,
\code{\_choice\_iterator(self, current, random\_gener\newline ator) ->
Iterator[\_SizingChoice]}, whose four branches \emph{are} the precise
definition of what distinguishes the four heuristics: \code{RANDOM\_GREEDY}
takes \code{\_critical\_path\_one\_step\_c\newline hoices(current)} and shuffles it
with \code{random\_generator}; \code{CRITICALITY\_EFFORT\_GAP} passes that
same one-step choice list through
\code{iter\_criticality\_effort\_gap\_choices}; \code{BRUTE\_FORCE} and
\code{SLACK\_WEIGHTED\_CAPACITANCE} both instead start from
\code{\_eligible\_choices(current, allowed\_sizes)}, flattened directly for
brute force or passed through \code{iter\_ranked\_optimi\newline zation\_choices}
for the canonical heuristic.

Two \emph{sources} of candidates are used, and it is important not to
conflate them:

\begin{itemize}
    \item \code{\_critical\_path\_one\_step\_choices(current) ->
          list[tuple[Gate, Cell]]} returns, for every gate on any current
          critical path, only \emph{one} candidate: its immediate
          next-larger cell in the same family (never a logical-effort-derived
          bracket, never a jump of more than one discrete size step).
          \textbf{Both} \code{RANDOM\_GREEDY} \textbf{and}
          \code{CRITICALITY\_EFFORT\_GAP} draw from this same, narrower
          candidate set --- they differ from each other only in the
          \emph{order} they visit it (uniformly shuffled versus
          score-ranked), not in which cells are candidates at all.
    \item \code{\_eligible\_choices(current, allowed\_sizes) ->
          list[OptimizationChoice]} ( where \code{OptimizationChoice =
          tuple[Gate, tuple[Cell, ...]]}) calls
          \code{analysis/ logica\_effort.py:get\_optimization\_candidates},
          which returns, for every gate on any current critical path,
          \emph{every} discrete cell in its logical-effort sizing bracket
          (\cref{sec:le-candidates}) --- potentially more than one size
          step in a single proposal. \textbf{Both} \code{BRUTE\_FORCE}
          \textbf{and} \code{SLACK\_WEIGHTED\_CAPACITANCE} draw from this
          wider, logical-effort-derived candidate set.
\end{itemize}

\subsection{Brute Force}
\label{sec:opt-bruteforce}

Every $(gate, candidate)$ pair in the logical-effort bracket set is
evaluated, in full, every iteration, and \code{\_select\_next\_change}
picks the single best. This is the most thorough of the four heuristics
--- and, per \cref{sec:opt-accept}, the only one that genuinely compares
multiple candidates against each other before choosing, rather than
walking a pre-ordered list --- at the cost of the largest number of full
STA re-runs per accepted change. It is available for manual and diagnostic
use but is not the heuristic invoked to produce the canonical optimized
result reported in \cref{ch:results}.

\subsection{Slack-Weighted Capacitance (Canonical)}
\label{sec:opt-swc}

Draws from the same logical-effort bracket set as brute force, but visits
it as a sequential, first-improvement search, in an order fixed once per
outer-loop iteration by \code{iter\_ranked\_op\newline timization\_choices(circuit,
paths, choices) -> Iterator[tuple[Gate, Cell]]},\newline which sorts candidates by
\code{(-score, gate\_position)} --- highest score first, ties broken by
netlist declaration order for determinism --- where the score is computed
by the private helper \code{\_slack\_weighted\_capacitance\_scores(circuit,
paths) -> dict[str, float]}: for every stage of every current critical
path, an \emph{urgency} factor (ranging from $1.0$, for the
least-violating path present, up to $2.0$, for the most-violating one,
linearly interpolated by that path's slack relative to the range of
slacks across all reported critical paths) is multiplied by the
logarithmic mismatch $|\ln(C^{*}_{\mathrm{in},i} / C_{\mathrm{in},i})|$
between the stage's logical-effort target capacitance
(\cref{eq:le-target-cap}) and its currently-assigned capacitance, and the
contributions are \emph{summed} across every path a given gate appears on.
This is the heuristic actually invoked by the standard analysis pipeline
(\code{CANONICAL\_OPTIMIZATION} in \code{app/experiments.py}) to produce
the reported ``logical-effort-guided'' result throughout \cref{ch:results}.

\subsection{Criticality/Effort-Gap}
\label{sec:opt-ceg}

A second logical-effort-\emph{informed} heuristic, but one that --- unlike
slack-weighted capacitance --- does \emph{not} draw from the capacitance
bracket at all. It re-orders the same narrow, one-step-up candidate set
used by random greedy (\code{\_critical\_path\_one\_step\_choices}) via
\code{iter\_criticality\_effort\_gap\_choices(circuit, paths, choices)},
whose ranking is produced by \code{criticality\_effort\_gap\_scores(circuit,
paths) -> dict[str, float]} (documented in-source as scoring gates by
``urgency \(\times\) (current stage effort $-$ optimal effort)''): the same
urgency factor as \cref{sec:opt-swc}, but multiplied instead by the raw
logical-effort stage-effort overshoot $f_i - \hat f$
(\cref{eq:le-min-delay}) directly --- and, where slack-weighted capacitance
\emph{sums} a gate's contribution across every path it appears on, this
score takes the \emph{maximum} over paths instead. In short: slack-weighted
capacitance asks ``how large a capacitance mismatch does this gate have,
aggregated over everywhere it's used, and how large a library jump would
fix it,'' while criticality/effort-gap asks ``how far over its ideal
per-stage effort is this gate, at its single worst appearance, if I only
allow myself the next library size up.'' Because both heuristics can, in
principle, still propose the exact same one-step-up resize when that
happens to also lie in the logical-effort bracket, a difference in final
outcome between the two is not guaranteed to come from candidate
\emph{scope} at all --- \cref{sec:res-case-divergence} traces one benchmark
circuit where it instead comes from \emph{ordering interacting with the
area/power budget}: criticality/effort-gap accepts an earlier, WNS-neutral
change that consumes part of the budget, after which the one candidate
that would have fixed the violation is rejected purely on that
now-tighter budget, whereas slack-weighted capacitance reaches the
identical candidate earlier, before any budget had been spent.

\subsection{Random Greedy (Baseline)}
\label{sec:opt-random-greedy}

The baseline heuristic against which the three logical-effort-informed
heuristics are compared: the same one-step-up candidate set as
criticality/effort-gap, visited in a uniformly random order
(\code{random\_generator.shuffle(choices)}, seeded by
\code{CircuitOptimizer(..., random\_seed=...)}) rather than any
score-based order. It is, in the strict sense, the only one of the four
heuristics that makes no use of logical effort anywhere in its candidate
generation.

\begin{table}[H]
\centering
\small
\begin{tabular}{@{}lccc@{}}
\toprule
\textbf{Heuristic} & \textbf{Candidate source} & \textbf{Ordering} & \textbf{Evaluation} \\
\midrule
Brute Force & LE bracket (\code{\_eligible\_choices}) & --- (all at once) & Best-of-batch \\
Slack-Wtd.\ Capacitance & LE bracket (\code{\_eligible\_choices}) & Score-ranked & Sequential \\
Criticality/Effort-Gap & One-step-up only & Score-ranked & Sequential \\
Random Greedy & One-step-up only & Uniform random & Sequential \\
\bottomrule
\end{tabular}
\caption{Summary of the four gate-sizing heuristics, by candidate source
and search strategy (\cref{sec:opt-heuristics}). All four share the
identical two-phase accept/reject policy (\cref{sec:opt-accept}).}
\label{tab:heuristic-summary}
\end{table}
